\section{Reporting and ethics}
\label{sec:reporting-ethics}

The practical proposal is a reporting distinction that follows the role. Participant studies should report recruitment, eligibility, consent, relevant demographics, compensation, exclusion, and the statistical treatment of responses. Evaluation studies should report a different set \citep{aera2014standards}:

\begin{itemize}
  \item the evaluator's qualifications for the standard applied;
  \item the evaluation task and the criterion or diagnostic used;
  \item the materials and context available to the evaluator;
  \item whether the evaluator was blind to the hypothesis;
  \item independence from the author and the project;
  \item how disagreements were adjudicated;
  \item any conflict of interest or relationship to the work;
  \item whether the judgment concerns form, meaning, register, dialect, processing difficulty, or theoretical classification.
\end{itemize}

This routing does not weaken ethics review; it aims it. A study that targets persons, whether their processing, dialect, attitudes, or the social distribution of their judgments, belongs in participant review with the usual protections. A study that uses expert competence to assess linguistic materials belongs under the standards that govern measurement: calibration, reliability, independence, and conflict management \citep{aera2014standards}.

The institutional claim stays narrow. Local boards decide the scope of their own review, and no definition settles policy. The conceptual point, argued in \autoref{sec:regulatory-asymmetry}, is only that the word \enquote{judgment} should not trigger participant review on its own. What should trigger it is whether the design studies persons as persons.
